%! Equivalent Representations of Polynomial and Rational Expressions — Unit 1, Lesson 1.11 — Homework (Answer Key)
\documentclass[10pt]{article}
\usepackage{precalculus-article}
\usepackage{precalculus-key}   % pulls in -boxes; do NOT also load -boxes
\newcommand{\TallMath}[1]{$\displaystyle #1\rule[-1.4em]{0pt}{3.2em}$}

\begin{document}

\pageheader{Unit 1, Lesson 1.11}{Homework --- Answer Key}
\namedateperiod

\vspace{0.12in}

\begin{enumerate}[itemsep=18pt]

\item For each rational function, state which form (\textbf{factored} or \textbf{standard}) is
      more useful for the requested task, then answer the question.

      \begin{enumerate}[label=(\alph*), itemsep=10pt]
        \item $r(x) = \dfrac{x^2+3x-10}{x^2-4}$. \quad Find all zeros, holes, and VAs.

              \vspace{0.06in}
              Best form: \ans{factored}

              Factor numerator: \ans{$(x+5)(x-2)$} \quad Factor denominator: \ans{$(x-2)(x+2)$}

              Zeros: \ans{$x=-5$} \quad Holes: \ans{$\left(2,\,\tfrac{7}{4}\right)$} \quad VAs: \ans{$x=-2$}

        \item $s(x) = \dfrac{3x^5 - 2x^2 + 1}{4x^3 - x}$. \quad Describe the end behavior.

              \vspace{0.06in}
              Best form: \ans{standard}

              As $x \to +\infty$: $s(x) \to$ \ans{$+\infty$} \quad As $x \to -\infty$: $s(x) \to$ \ans{$+\infty$}
      \end{enumerate}

\item Perform polynomial long division. Show all steps.

      Divide $f(x) = 3x^3 - 5x^2 + 2x - 4$ by $g(x) = x - 2$.

      \vspace{0.1in}
      Work space: $3x^3\div x = 3x^2$; $3x^2(x-2)=3x^3-6x^2$; subtract $\to x^2+2x-4$;
      $x^2\div x = x$; $x(x-2)=x^2-2x$; subtract $\to 4x-4$;
      $4x\div x = 4$; $4(x-2)=4x-8$; subtract $\to 4$.

      $q(x) = $ \ans{$3x^2 + x + 4$} \quad $r = $ \ans{$4$}

      Write as $f(x) = g(x)\cdot q(x) + r$: \ans{$(x-2)(3x^2+x+4)+4$}

\item A rational function is defined by $\displaystyle R(x) = \dfrac{x^2 + x - 6}{x - 1}$.

      \begin{enumerate}[label=(\alph*), itemsep=8pt]
        \item Perform polynomial long division.

              \vspace{0.06in}
              $R(x) = $ \ans{$x+2+\dfrac{-4}{x-1}$}

        \item State the equation of the slant asymptote. Explain in one sentence why this is
              the slant asymptote.

              \vspace{0.06in}
              Slant asymptote: $y = $ \ans{$x+2$}

              \ansline{As $x\to\pm\infty$, the remainder term $\frac{-4}{x-1}\to 0$, so $R(x)$ approaches the quotient line $y=x+2$.}

        \item What is the remainder when $x^2+x-6$ is divided by $x-1$? What does this tell
              you about $R(1)$?

              \vspace{0.06in}
              Remainder: \ans{$-4$} \quad $R(1) = $ \ans{undefined ($x=1$ is not in the domain)}
      \end{enumerate}

\item Use Pascal's Triangle to expand each expression. Show the row of Pascal's Triangle you used.

      \begin{enumerate}[label=(\alph*), itemsep=10pt]
        \item $(x + 3)^3$

              \vspace{0.06in}
              Row used: \ans{Row 3: $1\quad 3\quad 3\quad 1$}

              $(x+3)^3 = $ \ans{$x^3 + 9x^2 + 27x + 27$}

        \item $(2x - 1)^4$ \quad \textit{(Hint: let $a = 2x$ and $b = -1$)}

              \vspace{0.06in}
              Row used: \ans{Row 4: $1\quad 4\quad 6\quad 4\quad 1$}

              $(2x-1)^4 = $ \ans{$16x^4 - 32x^3 + 24x^2 - 8x + 1$}
      \end{enumerate}

\item \textbf{(AP-Style Multiple Choice)} Which of the following statements about
      $\displaystyle p(x) = \dfrac{x^3 - 8}{x - 2}$ is correct?

      \medskip
      \begin{enumerate}[label=(\Alph*), itemsep=4pt]
        \item $p$ has a vertical asymptote at $x = 2$.
        \item \ans{$p$ has a hole at $x = 2$ with $y$-coordinate $12$.} \textcolor{keyred}{\textbf{$\leftarrow$ correct}}
        \item $p$ simplifies to $x^2 - 2x + 4$ for all real $x$.
        \item $p$ has a slant asymptote of $y = x^2 + 2x + 4$.
      \end{enumerate}

      \vspace{0.06in}
      Answer: \ans{B} \quad Show your reasoning:

      \vspace{0.06in}
      \ansline{Factor: $x^3-8 = (x-2)(x^2+2x+4)$. Shared factor $(x-2)$ in numerator and denominator with equal multiplicity, so $x=2$ is a hole. Substitute $x=2$ into $x^2+2x+4$: $4+4+4=12$. Hole at $(2,12)$.}

\end{enumerate}

\begin{teachernote}
1(a): Hole at $x=2$: substitute into simplified $\frac{x+5}{x+2}$ to get $\frac{7}{4}$.

1(b): Leading terms $3x^5$ and $4x^3$ give ratio $\frac{3}{4}x^2 \to +\infty$ as $x\to\pm\infty$, but signs differ. As $x\to+\infty$: $+\infty$; as $x\to-\infty$: $-\infty$ (odd degree in quotient).

Wait --- $\frac{3x^5}{4x^3} = \frac{3}{4}x^2$ which goes to $+\infty$ in both directions. Correction: as $x\to-\infty$, $s(x)\to+\infty$ as well (even power). Teacher should verify this with the class and note the answer key reflects $+\infty$ for both directions.

4(b): $(2x)^4 = 16x^4$; $4(2x)^3(-1) = -32x^3$; $6(2x)^2(1)=24x^2$; $4(2x)(-1)^3=-8x$; $(-1)^4=1$.
\end{teachernote}

\end{document}
